\documentclass[12pt, ukrainian]{article}
\usepackage[a4paper]{geometry}
\setlength\parindent{0mm}
\geometry{top=1.1cm,bottom=1.1cm,left=1.1cm,right=1.1cm}
\pagestyle{empty}
\usepackage[T2A]{fontenc}
\usepackage[utf8]{inputenc}
\usepackage[ukrainian]{babel}
\usepackage{graphicx}
\usepackage{caption}
\DeclareCaptionFormat{nocaption}{}
\captionsetup{format=nocaption,aboveskip=0pt,belowskip=0pt}
\usepackage[Export]{adjustbox} 
\adjustboxset{max size={0.9\linewidth}{0.9\paperheight}}
\begin{document}
%begin
\begin {Large}
%beginitem
1. Написати функцію, що для інтервалу $[a, b)$ знаходить найближчу пару чисел з цього інтервалу, що мають не менше трьох простих дільників та мають цифру 3 у своєму десятковому запису.\\	Якість коду також оцінюється. 



\vspace{10mm}

%enditem
\newpage

%beginitem
2. 
Написати функцію, що виводить послідовність \underline{цілих} чисел
$$\scalebox{1.2}{$(-1)^{k}C_{n+k}^{2k-1}$}, \quad k=1,\ldots,n$$
Оцінюється економність обчислень. Використовувати ** та виклики функцій заборонено.

%enditem
\end {Large}\newpage

%end
%end



%begin
\begin {Large}
%beginitem
1. Написати функцію, що для інтервалу $[a, b)$ знаходить знаходить найбільший степінь числа 3, на який ділиться хоча б одне ціле число з цього інтервалу.\\	Якість коду також оцінюється. 



\vspace{10mm}

%enditem
\newpage

%beginitem
2. 
Написати функцію, що виводить послідовність \underline{цілих} чисел
$$\scalebox{1.2}{$(-1)^{k}C_{n+k}^{k+1}$}, \quad k=0,1,\ldots,n$$
Оцінюється економність обчислень. Використовувати ** та виклики функцій заборонено.

%enditem
\end {Large}\newpage

%end
%end



%begin
\begin {Large}
%beginitem
1. Написати функцію, що для інтервалу $[a, b)$ знаходить найближчу пару чисел з цього інтервалу, що мають рівно два простих дільники.\\	Якість коду також оцінюється. 



\vspace{10mm}

%enditem
\newpage

%beginitem
2. 
Написати функцію, що виводить послідовність \underline{цілих} чисел
$$\scalebox{1.2}{$\frac{C_{2k}^k}{k+1}$}, \quad k=0,1,\ldots,n$$
Оцінюється економність обчислень. Використовувати ** та виклики функцій заборонено.

%enditem
\end {Large}\newpage

%end
%end



%begin
\begin {Large}
%beginitem
1. Написати функцію, що для інтервалу $[a, b)$ знаходить найближчу пару чисел з цього інтервалу, що мають парну кількість різних простих дільників. Якщо таких пар кілька, то цікавить пара з найменшим добутком.\\	Якість коду також оцінюється. 



\vspace{10mm}

%enditem
\newpage

%beginitem
2. 
Написати функцію, що виводить послідовність \underline{цілих} чисел
$$\scalebox{1.2}{$C_{n+k}^{2k}$}, \quad k=1,\ldots,n$$
Оцінюється економність обчислень. Використовувати ** та виклики функцій заборонено.

%enditem
\end {Large}\newpage

%end
%end



%begin
\begin {Large}
%beginitem
1. Написати функцію, що в інтервалі $[a, b)$ знаходить ціле число, що ділиться на найбільший степінь числа 3. Якщо таких чисел декілька, повернути найбільше.\\	Якість коду також оцінюється. 



\vspace{10mm}

%enditem
\newpage

%beginitem
2. 
Написати функцію, що виводить послідовність \underline{цілих} чисел
$$\scalebox{1.2}{$(-1)^{k}\frac{(2k+1)!!}{(k!)^2}$}, \quad k=1,\ldots,n$$
Оцінюється економність обчислень. Використовувати ** та виклики функцій заборонено.

%enditem
\end {Large}\newpage

%end
%end



%begin
\begin {Large}
%beginitem
1. Написати функцію, що для інтервалу $[a, b)$ знаходить найближчу пару чисел з цього інтервалу, що мають непарну кількість різних простих дільників. Якщо таких пар кілька, то цікавить пара з найменшим добутком.\\	Якість коду також оцінюється. 



\vspace{10mm}

%enditem
\newpage

%beginitem
2. 
Написати функцію, що виводить послідовність \underline{цілих} чисел
$$\scalebox{1.2}{$C_{n+k}^{2k-1}$}, \quad k=1,\ldots,n$$
Оцінюється економність обчислень. Використовувати ** та виклики функцій заборонено.

%enditem
\end {Large}\newpage

%end
%end



%begin
\begin {Large}
%beginitem
1. Написати функцію, що для інтервалу $[a, b)$ знаходить найближчу пару чисел з цього інтервалу, що мають парну кількість різних простих дільників. Якщо таких пар кілька, то цікавить пара з найбільшою сумою.\\	Якість коду також оцінюється. 



\vspace{10mm}

%enditem
\newpage

%beginitem
2. 
Написати функцію, що виводить послідовність \underline{цілих} чисел
$$\scalebox{1.2}{$(-1)^{k}C_{2n+k}^{n+k}$}, \quad k=1,\ldots,n$$
Оцінюється економність обчислень. Використовувати ** та виклики функцій заборонено.

%enditem
\end {Large}\newpage

%end
%end



%begin
\begin {Large}
%beginitem
1. Написати функцію, що в інтервалі $[a, b)$ знаходить ціле число, що ділиться на найбільший степінь числа 3. Якщо таких чисел декілька, повернути найменше.\\	Якість коду також оцінюється. 



\vspace{10mm}

%enditem
\newpage

%beginitem
2. 
Написати функцію, що виводить послідовність \underline{цілих} чисел
$$\scalebox{1.2}{$(-1)^{k}C_{n+k}^{k}$}, \quad k=1,\ldots,n$$
Оцінюється економність обчислень. Використовувати ** та виклики функцій заборонено.

%enditem
\end {Large}\newpage

%end
%end



%begin
\begin {Large}
%beginitem
1. Написати функцію, що для інтервалу $[a, b)$ знаходить найближчу пару чисел з цього інтервалу, що мають непарну кількість різних простих дільників. Якщо таких пар кілька, то цікавить пара з найбільшою сумою.\\	Якість коду також оцінюється. 



\vspace{10mm}

%enditem
\newpage

%beginitem
2. 
Написати функцію, що виводить послідовність \underline{цілих} чисел
$$\scalebox{1.2}{$(-1)^{k}C_{n+k}^{2k}$}, \quad k=0,1,\ldots,n$$
Оцінюється економність обчислень. Використовувати ** та виклики функцій заборонено.

%enditem
\end {Large}\newpage

%end
%end



%begin
\begin {Large}
%beginitem
1. Написати функцію, що для інтервалу $[a, b)$ знаходить кількість чисел, квадрати яких діляться на суму квадратів їх цифр. (Розглядаємо десятковий запис.)\\	Якість коду також оцінюється. 



\vspace{10mm}

%enditem
\newpage

%beginitem
2. 
Написати функцію, що виводить послідовність \underline{цілих} чисел
$$\scalebox{1.2}{$C_{n+2k}^{k}$}, \quad k=0,1,\ldots,n$$
Оцінюється економність обчислень. Використовувати ** та виклики функцій заборонено.

%enditem
\end {Large}\newpage

%end
%end



%begin
\begin {Large}
%beginitem
1. Написати функцію, що в інтервалі $[a, b)$ знаходить ціле число, що ділиться на найбільший степінь числа 3. Якщо таких чисел декілька, повернути найменше.\\	Якість коду також оцінюється. 



\vspace{10mm}

%enditem
\newpage

%beginitem
2. 
Написати функцію, що виводить послідовність \underline{цілих} чисел
$$\scalebox{1.2}{$\frac{(2k+1)!!}{(k!)^2}$}, \quad k=1,\ldots,n$$
Оцінюється економність обчислень. Використовувати ** та виклики функцій заборонено.

%enditem
\end {Large}\newpage

%end
%end



%begin
\begin {Large}
%beginitem
1. Написати функцію, що для інтервалу $[a, b)$ знаходить кількість чисел, квадрати яких діляться на суму квадратів їх цифр. (Розглядаємо десятковий запис.)\\	Якість коду також оцінюється. 



\vspace{10mm}

%enditem
\newpage

%beginitem
2. 
Написати функцію, що виводить послідовність \underline{цілих} чисел
$$\scalebox{1.2}{$(-1)^{k}C_{n+k}^{k}$}, \quad k=0,1,\ldots,n$$
Оцінюється економність обчислень. Використовувати ** та виклики функцій заборонено.

%enditem
\end {Large}\newpage

%end
%end



%begin
\begin {Large}
%beginitem
1. Написати функцію, що для інтервалу $[a, b)$ знаходить найближчу пару чисел з цього інтервалу, що мають парну кількість різних простих дільників. Якщо таких пар кілька, то цікавить пара з найменшим добутком.\\	Якість коду також оцінюється. 



\vspace{10mm}

%enditem
\newpage

%beginitem
2. 
Написати функцію, що виводить послідовність \underline{цілих} чисел
$$\scalebox{1.2}{$C_{n+k}^{k+1}$}, \quad k=0,1,\ldots,n$$
Оцінюється економність обчислень. Використовувати ** та виклики функцій заборонено.

%enditem
\end {Large}\newpage

%end
%end



%begin
\begin {Large}
%beginitem
1. Написати функцію, що для інтервалу $[a, b)$ знаходить найближчу пару чисел з цього інтервалу, що мають рівно два простих дільники.\\	Якість коду також оцінюється. 



\vspace{10mm}

%enditem
\newpage

%beginitem
2. 
Написати функцію, що виводить послідовність \underline{цілих} чисел
$$\scalebox{1.2}{$(-1)^{k}\frac{(2k+1)!!}{(k+1)!(k-1)!}$}, \quad k=2,\ldots,n$$
Оцінюється економність обчислень. Використовувати ** та виклики функцій заборонено.

%enditem
\end {Large}\newpage

%end
%end



%begin
\begin {Large}
%beginitem
1. Написати функцію, що для інтервалу $[a, b)$ знаходить знаходить найбільший степінь числа 3, на який ділиться хоча б одне ціле число з цього інтервалу.\\	Якість коду також оцінюється. 



\vspace{10mm}

%enditem
\newpage

%beginitem
2. 
Написати функцію, що виводить послідовність \underline{цілих} чисел
$$\scalebox{1.2}{$C_{2n+k}^{n+k}$}, \quad k=1,\ldots,n$$
Оцінюється економність обчислень. Використовувати ** та виклики функцій заборонено.

%enditem
\end {Large}\newpage

%end
%end



%begin
\begin {Large}
%beginitem
1. Написати функцію, що в інтервалі $[a, b)$ знаходить ціле число, що ділиться на найбільший степінь числа 3. Якщо таких чисел декілька, повернути найбільше.\\	Якість коду також оцінюється. 



\vspace{10mm}

%enditem
\newpage

%beginitem
2. 
Написати функцію, що виводить послідовність \underline{цілих} чисел
$$\scalebox{1.2}{$(-1)^{k}\frac{(2k)!!}{(k+1)!(k-1)!}$}, \quad k=2,\ldots,n$$
Оцінюється економність обчислень. Використовувати ** та виклики функцій заборонено.

%enditem
\end {Large}\newpage

%end
%end



%begin
\begin {Large}
%beginitem
1. Написати функцію, що для інтервалу $[a, b)$ знаходить найближчу пару чисел з цього інтервалу, що мають непарну кількість різних простих дільників. Якщо таких пар кілька, то цікавить пара з найменшим добутком.\\	Якість коду також оцінюється. 



\vspace{10mm}

%enditem
\newpage

%beginitem
2. 
Написати функцію, що виводить послідовність \underline{цілих} чисел
$$\scalebox{1.2}{$(-1)^{k}C_{n+2k}^{k}$}, \quad k=0,1,\ldots,n$$
Оцінюється економність обчислень. Використовувати ** та виклики функцій заборонено.

%enditem
\end {Large}\newpage

%end
%end



%begin
\begin {Large}
%beginitem
1. Написати функцію, що для інтервалу $[a, b)$ знаходить найближчу пару чисел з цього інтервалу, що мають парну кількість різних простих дільників. Якщо таких пар кілька, то цікавить пара з найбільшою сумою.\\	Якість коду також оцінюється. 



\vspace{10mm}

%enditem
\newpage

%beginitem
2. 
Написати функцію, що виводить послідовність \underline{цілих} чисел
$$\scalebox{1.2}{$(-1)^{k}C_{2k+1}^{k}$}, \quad k=2,\ldots,n$$
Оцінюється економність обчислень. Використовувати ** та виклики функцій заборонено.

%enditem
\end {Large}\newpage

%end
%end



%begin
\begin {Large}
%beginitem
1. Написати функцію, що для інтервалу $[a, b)$ знаходить найближчу пару чисел з цього інтервалу, що мають не менше трьох простих дільників та мають цифру 3 у своєму десятковому запису.\\	Якість коду також оцінюється. 



\vspace{10mm}

%enditem
\newpage

%beginitem
2. 
Написати функцію, що виводить послідовність \underline{цілих} чисел
$$\scalebox{1.2}{$C_{2k+1}^{k}$}, \quad k=2,\ldots,n$$
Оцінюється економність обчислень. Використовувати ** та виклики функцій заборонено.

%enditem
\end {Large}\newpage

%end
%end



%begin
\begin {Large}
%beginitem
1. Написати функцію, що для інтервалу $[a, b)$ знаходить найближчу пару чисел з цього інтервалу, що мають непарну кількість різних простих дільників. Якщо таких пар кілька, то цікавить пара з найбільшою сумою.\\	Якість коду також оцінюється. 



\vspace{10mm}

%enditem
\newpage

%beginitem
2. 
Написати функцію, що виводить послідовність \underline{цілих} чисел
$$\scalebox{1.2}{$\frac{(2k)!!}{(k+1)!(k-1)!}$}, \quad k=2,\ldots,n$$
Оцінюється економність обчислень. Використовувати ** та виклики функцій заборонено.

%enditem
\end {Large}\newpage

%end
%end



%begin
\begin {Large}
%beginitem
1. Написати функцію, що для інтервалу $[a, b)$ знаходить найближчу пару чисел з цього інтервалу, що мають рівно два простих дільники.\\	Якість коду також оцінюється. 



\vspace{10mm}

%enditem
\newpage

%beginitem
2. 
Написати функцію, що виводить послідовність \underline{цілих} чисел
$$\scalebox{1.2}{$(-1)^{k}C_{2n+k+1}^{n+k}$}, \quad k=0,1,\ldots,n$$
Оцінюється економність обчислень. Використовувати ** та виклики функцій заборонено.

%enditem
\end {Large}\newpage

%end
%end



%begin
\begin {Large}
%beginitem
1. Написати функцію, що для інтервалу $[a, b)$ знаходить кількість чисел, квадрати яких діляться на суму квадратів їх цифр. (Розглядаємо десятковий запис.)\\	Якість коду також оцінюється. 



\vspace{10mm}

%enditem
\newpage

%beginitem
2. 
Написати функцію, що виводить послідовність \underline{цілих} чисел
$$\scalebox{1.2}{$(-1)^{k}C_{2n+k}^{n+k}$}, \quad k=0,1,\ldots,n$$
Оцінюється економність обчислень. Використовувати ** та виклики функцій заборонено.

%enditem
\end {Large}\newpage

%end
%end



%begin
\begin {Large}
%beginitem
1. Написати функцію, що в інтервалі $[a, b)$ знаходить ціле число, що ділиться на найбільший степінь числа 3. Якщо таких чисел декілька, повернути найменше.\\	Якість коду також оцінюється. 



\vspace{10mm}

%enditem
\newpage

%beginitem
2. 
Написати функцію, що виводить послідовність \underline{цілих} чисел
$$\scalebox{1.2}{$C_{2n+k}^{n+k}$}, \quad k=0,1,\ldots,n$$
Оцінюється економність обчислень. Використовувати ** та виклики функцій заборонено.

%enditem
\end {Large}\newpage

%end
%end



%begin
\begin {Large}
%beginitem
1. Написати функцію, що в інтервалі $[a, b)$ знаходить ціле число, що ділиться на найбільший степінь числа 3. Якщо таких чисел декілька, повернути найбільше.\\	Якість коду також оцінюється. 



\vspace{10mm}

%enditem
\newpage

%beginitem
2. 
Написати функцію, що виводить послідовність \underline{цілих} чисел
$$\scalebox{1.2}{$C_{n+k}^{k}$}, \quad k=0,1,\ldots,n$$
Оцінюється економність обчислень. Використовувати ** та виклики функцій заборонено.

%enditem
\end {Large}\newpage

%end
%end



%begin
\begin {Large}
%beginitem
1. Написати функцію, що для інтервалу $[a, b)$ знаходить найближчу пару чисел з цього інтервалу, що мають не менше трьох простих дільників та мають цифру 3 у своєму десятковому запису.\\	Якість коду також оцінюється. 



\vspace{10mm}

%enditem
\newpage

%beginitem
2. 
Написати функцію, що виводить послідовність \underline{цілих} чисел
$$\scalebox{1.2}{$(-1)^{k}C_{n+k}^{2k}$}, \quad k=1,\ldots,n$$
Оцінюється економність обчислень. Використовувати ** та виклики функцій заборонено.

%enditem
\end {Large}\newpage

%end
%end



%begin
\begin {Large}
%beginitem
1. Написати функцію, що для інтервалу $[a, b)$ знаходить найближчу пару чисел з цього інтервалу, що мають непарну кількість різних простих дільників. Якщо таких пар кілька, то цікавить пара з найменшим добутком.\\	Якість коду також оцінюється. 



\vspace{10mm}

%enditem
\newpage

%beginitem
2. 
Написати функцію, що виводить послідовність \underline{цілих} чисел
$$\scalebox{1.2}{$C_{2n+k+1}^{n+k}$}, \quad k=0,1,\ldots,n$$
Оцінюється економність обчислень. Використовувати ** та виклики функцій заборонено.

%enditem
\end {Large}\newpage

%end
%end



\end{document}

